\chapter[FORMATS COMPRESSÉS DE RAPPORT DE POSITION]{Formats compressés de rapport de
position}\label{formats-compressuxe9s-de-rapport-de-position}

En format compressé, le champ Information contient lat/long +
course/speed ou pré-calculated range ou altitude, comprimés pour
minimiser la longueur du paquet (meilleure chance de réception en
conditions dégradées). Symbol Code présent. Commentaire libre optionnel
non compressé.

\section*{Avantages}\label{avantages}

\begin{itemize}
\tightlist
\item
  Rétrocompatible avec les formats existants (\texttt{!}, \texttt{/},
  \texttt{@}, \texttt{=}).
\item
  Supporte n'importe quel commentaire.
\item
  Vitesse précise à ±1 mph jusqu'à 40 mph, 3 \% à 600 mph.
\item
  Altitude en pieds précise à ±0.4 \% de 1 ft à 3000 miles.
\item
  Algorithme unique pour lat/long compressées.
\item
  Position améliorée à 1 pied dans le monde entier.
\item
  Pre-calculated range sur 1 octet.
\item
  Potentiel de compression 50 \% de tout format de position.
\item
  Réduction potentielle de 40 \% de la longueur du NMEA brut.
\item
  Réduction supplémentaire de 7 octets pour les altitudes GGA.
\item
  Support de la compression TNC à la source NMEA.
\item
  Compression digipeater à la volée des anciens trackers NMEA.
\item
  Toujours optionnel.
\end{itemize}

Inconvénients mineurs : le cap ne résout qu'à ±2°, et pas de support
PHG.

\section*{Format des données
compressées}\label{format-des-donnuxe9es-compressuxe9es}

Champ fixe 13 caractères : \texttt{/YYYYXXXX\$csT} où : - \texttt{/} =
Symbol Table Identifier - \texttt{YYYY} = latitude compressée -
\texttt{XXXX} = longitude compressée - \texttt{\$} = Symbol Code -
\texttt{cs} = course/speed OU range pré-calculé OU altitude compressée -
\texttt{T} = indicateur de type de compression

\begin{longtable}[]{@{}
  >{\raggedright\arraybackslash}p{(\columnwidth - 12\tabcolsep) * \real{0.1429}}
  >{\raggedright\arraybackslash}p{(\columnwidth - 12\tabcolsep) * \real{0.1429}}
  >{\raggedright\arraybackslash}p{(\columnwidth - 12\tabcolsep) * \real{0.1429}}
  >{\raggedright\arraybackslash}p{(\columnwidth - 12\tabcolsep) * \real{0.1429}}
  >{\raggedright\arraybackslash}p{(\columnwidth - 12\tabcolsep) * \real{0.1429}}
  >{\raggedright\arraybackslash}p{(\columnwidth - 12\tabcolsep) * \real{0.1429}}
  >{\raggedright\arraybackslash}p{(\columnwidth - 12\tabcolsep) * \real{0.1429}}@{}}
\toprule\noalign{}
\begin{minipage}[b]{\linewidth}\raggedright
Champ
\end{minipage} & \begin{minipage}[b]{\linewidth}\raggedright
Sym Tbl ID
\end{minipage} & \begin{minipage}[b]{\linewidth}\raggedright
Comp Lat (YYYY)
\end{minipage} & \begin{minipage}[b]{\linewidth}\raggedright
Comp Long (XXXX)
\end{minipage} & \begin{minipage}[b]{\linewidth}\raggedright
Symbol Code
\end{minipage} & \begin{minipage}[b]{\linewidth}\raggedright
cs
\end{minipage} & \begin{minipage}[b]{\linewidth}\raggedright
Comp Type (T)
\end{minipage} \\
\midrule\noalign{}
\endhead
\bottomrule\noalign{}
\endlastfoot
Octets & 1 & 4 & 4 & 1 & 2 & 1 \\
\end{longtable}

Tous les octets sauf \texttt{/} et le Symbol Code sont ASCII imprimables
base-91 (\texttt{!..\{}), convertis en numérique par soustraction de 33.
Ex : \texttt{\#} a code ASCII 35, valeur numérique 2.

\section*{Symbole}\label{symbole}

Le Symbol Table Identifier en tête (au lieu d'un chiffre) indique un
rapport compressé, pas un lat/long normal.

\section*{Encodage Lat/Long}\label{encodage-latlong}

\begin{verbatim}
YYYY = 380926 × (90 - latitude)      [base 91]
XXXX = 190463 × (180 + longitude)    [base 91]
\end{verbatim}

Latitude positive au nord, longitude positive à l'est (degrés décimaux).

\textbf{Exemple} longitude -72.75° :

\begin{verbatim}
190463 × (180 - 72.75) = 20427156
20427156 / 91³ = 27, reste 80739
   80739 / 91² =  9, reste 6210
    6210 / 91¹ = 68, reste 22
\end{verbatim}

Ajout 33 : 60, 42, 101, 55 → \texttt{\textless{}*e7}.

\section*{Décodage Lat/Long}\label{duxe9codage-latlong}

Pour \texttt{YYYY\ =\ y1y2y3y4} et \texttt{XXXX\ =\ x1x2x3x4} :

\begin{verbatim}
Lat  =  90 - ((y1-33)×91³ + (y2-33)×91² + (y3-33)×91 + (y4-33)) / 380926
Long = -180 + ((x1-33)×91³ + (x2-33)×91² + (x3-33)×91 + (x4-33)) / 190463
\end{verbatim}

Ex \texttt{\textless{}*e7} :

\begin{verbatim}
Long = -180 + (27×91³ + 9×91² + 68×91 + 22) / 190463
     = -180 + (20346417 + 74529 + 6188 + 22) / 190463
     = -180 + 107.25
     = -72.75°
\end{verbatim}

\section*{Course/Speed, Range,
Altitude}\label{coursespeed-range-altitude}

Les 2 octets \texttt{cs} après le Symbol Code contiennent, en base 91 :
- soit course/speed - soit range pré-calculé - soit altitude

Cas spécial \texttt{c\ =\ V} (espace) : pas de course/speed/range →
\texttt{csT} ignoré.

\textbf{Course/Speed} --- si \texttt{c} ASCII est dans \texttt{!..z}
(0-89 après -33) :

\begin{verbatim}
course = c × 4
speed  = 1.08^s - 1
\end{verbatim}

Ex \texttt{7P} : c=22, s=47 → course=88°, speed = 1.08⁴⁷ - 1 ≈ 36.2
knots.

\textbf{Pre-calculated Range} --- si \texttt{c\ =\ \{} :

\begin{verbatim}
range = 2 × 1.08^s
\end{verbatim}

Ex \texttt{\{?} : s=30 → range = 2 × 1.08³⁰ ≈ 20 miles.

\section*{Octet Compression Type (T)}\label{octet-compression-type-t}

Suit \texttt{cs}. Contient des bit-fields (GPS fix, source NMEA, origine
compression). Non significatif si \texttt{c\ =\ V}.

\begin{longtable}[]{@{}
  >{\raggedright\arraybackslash}p{(\columnwidth - 16\tabcolsep) * \real{0.1111}}
  >{\raggedright\arraybackslash}p{(\columnwidth - 16\tabcolsep) * \real{0.1111}}
  >{\raggedright\arraybackslash}p{(\columnwidth - 16\tabcolsep) * \real{0.1111}}
  >{\raggedright\arraybackslash}p{(\columnwidth - 16\tabcolsep) * \real{0.1111}}
  >{\raggedright\arraybackslash}p{(\columnwidth - 16\tabcolsep) * \real{0.1111}}
  >{\raggedright\arraybackslash}p{(\columnwidth - 16\tabcolsep) * \real{0.1111}}
  >{\raggedright\arraybackslash}p{(\columnwidth - 16\tabcolsep) * \real{0.1111}}
  >{\raggedright\arraybackslash}p{(\columnwidth - 16\tabcolsep) * \real{0.1111}}
  >{\raggedright\arraybackslash}p{(\columnwidth - 16\tabcolsep) * \real{0.1111}}@{}}
\toprule\noalign{}
\begin{minipage}[b]{\linewidth}\raggedright
Bit
\end{minipage} & \begin{minipage}[b]{\linewidth}\raggedright
7
\end{minipage} & \begin{minipage}[b]{\linewidth}\raggedright
6
\end{minipage} & \begin{minipage}[b]{\linewidth}\raggedright
5
\end{minipage} & \begin{minipage}[b]{\linewidth}\raggedright
4
\end{minipage} & \begin{minipage}[b]{\linewidth}\raggedright
3
\end{minipage} & \begin{minipage}[b]{\linewidth}\raggedright
2
\end{minipage} & \begin{minipage}[b]{\linewidth}\raggedright
1
\end{minipage} & \begin{minipage}[b]{\linewidth}\raggedright
0
\end{minipage} \\
\midrule\noalign{}
\endhead
\bottomrule\noalign{}
\endlastfoot
Valeur & non utilisé & non utilisé & GPS Fix & NMEA Source (2 bits) & &
Compression Origin (3 bits) & & \\
\end{longtable}

\begin{itemize}
\tightlist
\item
  GPS Fix : 0 = ancien, 1 = courant.
\item
  NMEA Source : 00 = other, 01 = GLL, 10 = GGA, 11 = RMC.
\item
  Compression Origin : 000 = Compressed, 001 = TNC BText, 010 = Software
  (DOS/Mac/Win/+SA), 011 = tbd, 100 = KPC3, 101 = Pico, 110 = Other
  tracker, 111 = Digipeater conversion.
\end{itemize}

\textbf{Exemple} : RMC + fix courant + APRSdos → binaire
\texttt{0\ 0\ 1\ 11\ 010} = 58 décimal. +33 = 91 = \texttt{{[}}.

Complet avec les exemples plus haut (49° 30' N, 72° 45' W, 36.2 knots,
88°, symbole « car »):

\begin{verbatim}
/ YYYY  XXXX  $  csT
/ 5L!!  <*e7  >  7P[
\end{verbatim}

\section*{Altitude (source GGA)}\label{altitude-source-gga}

Si le T byte indique GGA (bits 4-3 = \texttt{10}), la sentence contient
une altitude en pieds. Après compression, dans les octets \texttt{cs} :

\begin{verbatim}
altitude = 1.002^cs  feet
\end{verbatim}

où \texttt{cs} (numérique) = (c-33) × 91 + (s-33).

Ex \texttt{S{]}} : c=50, s=60 → cs = 50 × 91 + 60 = 4610 → altitude =
1.002⁴⁶¹⁰ = 10004 ft.

\section*{New Trackers}\label{new-trackers}

Firmware tracker peut compresser directement le GPS au format APRS
compressé. DTI \texttt{!} = temps réel + tracker sans capacité APRS. DTI
\texttt{/} = pas temps réel, avec dernier fix time inclus.

\section*{Old Trackers}\label{old-trackers}

Certains digipeaters convertissent le NMEA brut d'anciens trackers en
format compressé. Ils compressent si l'adresse de destination est
\texttt{GPS} (les 3 lettres exactes, pas \texttt{GPS*}). Les trackers ne
voulant pas être modifiés utilisent une autre destination générique
valide.

\section*{Formats de rapport
compressés}\label{formats-de-rapport-compressuxe9s}

\textbf{Sans timestamp} :

\begin{center}
\begin{tabular}{|c|c|c|c|c|c|c|c|}
\hline
\textbf{Champ} & \lit{!} ou \lit{=} & Sym Tbl ID & Comp Lat & Comp Long & Symbol Code & cs & T \\
\hline
\textbf{Octets} & 1 & 1 & 4 & 4 & 1 & 2 & 1 \\
\hline
\end{tabular}
\end{center}
\noindent + Comment (0-40 octets)

Exemples : - \texttt{=/5L!!\textless{}*e7\textgreater{}V\ sTComment} ---
messagerie, \texttt{V} espace après Symbol Code = pas de
cs/range/altitude, \texttt{sT} sont des fillers ignorés. -
\texttt{=/5L!!\textless{}*e7\textgreater{}7P{[}} --- messagerie, RMC,
avec course/speed. - \texttt{=/5L!!\textless{}*e7\textgreater{}\{?!} ---
messagerie, avec range. - \texttt{=/5L!!\textless{}*e7OS{]}S} ---
messagerie, GGA, altitude.

\textbf{Avec timestamp} :

\begin{center}
\begin{tabular}{|c|c|c|c|c|c|c|c|c|}
\hline
\textbf{Champ} & \lit{/} ou \lit{@} & Time & Sym Tbl ID & Comp Lat & Comp Long & Symbol Code & cs & T \\
\hline
\textbf{Octets} & 1 & 7 & 1 & 4 & 4 & 1 & 2 & 1 \\
\hline
\end{tabular}
\end{center}
\noindent + Comment (0-40 octets)

Ex : \texttt{@092345z/5L!!\textless{}*e7\textgreater{}\{?!} ---
messagerie, timestamp, range.

\begin{center}\rule{0.5\linewidth}{0.5pt}\end{center}
